%****************************************************************************************************
% Slide trình bày đề cương (khoảng 30 phút) — template Michaillat (MIT), biên dịch XeLaTeX
% xelatex slides && xelatex slides
% Quy ước: tiêu đề slide nêu kết luận; mọi số liệu lấy từ đề cương v14; số minh họa ghi rõ "minh họa".
%****************************************************************************************************
\documentclass[11pt, xcolor={dvipsnames}, aspectratio=169, usepdftitle=false]{beamer}
\usepackage{presentation}
\usepackage{tikz}
\usetikzlibrary{arrows.meta, positioning, calc}
\usepackage{tabularx}
\hypersetup{pdfpagemode=UseNone, hidelinks, pdfdisplaydoctitle=true,
  pdftitle={Hợp nhất sàn chứng khoán và chất lượng thị trường}}

% Ký hiệu
\newcommand{\HH}[1]{H\textsubscript{#1}}
\newcommand{\MM}[1]{M\textsubscript{#1}}
\newcommand{\vn}[1]{\mathrm{#1}}
\newcommand{\Post}{\vn{Post}}
\newcommand{\MQI}{\vn{MQI}}
\newcommand{\SMD}{\vn{SMD}}
% Nhấn mạnh dùng \al{} của template (một màu nhấn duy nhất)
\definecolor{cbBlue}{HTML}{0173B2}
\definecolor{cbOrange}{HTML}{DE8F05}
\newcolumntype{L}{>{\raggedright\arraybackslash}X}
\tikzset{box/.style={draw=black!70, rounded corners=2pt, align=center, font=\small, inner sep=5pt},
         arr/.style={-{Stealth[length=2.2mm]}, thick, black!70},
         darr/.style={-{Stealth[length=2.2mm]}, thick, dashed, black!70}}

\title{Hợp nhất sàn chứng khoán và chất lượng thị trường:\\ Bằng chứng thực nghiệm từ HNX và vị thế tài chính Hà Nội}
\author{Nguyễn Văn Trung\\[1mm]{\normalsize Đề cương nghiên cứu sinh, ngành Kinh tế và Quản lý, Trường Quốc tế, ĐHQGHN}}
\date{Hà Nội, 2026}

\begin{document}
\frame{\titlepage}

%=================================================================== MỞ ĐẦU
\begin{frame}
\frametitle{Ngày 28/12/2026, toàn bộ cổ phiếu niêm yết HNX chuyển sang HOSE}
\begin{center}
\begin{tikzpicture}[x=0.96cm]
\draw[arr] (0,0) -- (12.6,0) node[right, font=\small]{thời gian};
\foreach \x/\a/\b in {0.4/2021/TT 57: hạn năm 2023, 3.2/2023/TT 69: lùi đến hết 2025, 6.0/2025/TT 139: lùi đến 31/12/2026, 8.8/30/9/2026/VNX công bố lịch}
  {\draw[thick] (\x,0.12) -- (\x,-0.12); \node[above=4pt, font=\small\sbseries] at (\x,0.1){\a}; \node[below=4pt, font=\footnotesize, text width=2.2cm, align=center] at (\x,-0.1){\b};}
\draw[very thick, RubineRed] (11.3,0.25) -- (11.3,-0.25);
\node[above=4pt, font=\small\sbseries, RubineRed] at (11.3,0.2){28/12/2026};
\node[below=4pt, font=\footnotesize, text width=2.4cm, align=center, RubineRed] at (11.3,-0.2){phiên đầu tại HOSE};
\end{tikzpicture}
\end{center}
\vspace{2mm}
\begin{itemize}
\item \al{299 mã}, chuyển \al{cùng một ngày}, theo quyết định hành chính (doanh nghiệp không tự chọn)
\item Phiên cuối ở HNX 23/12; tạm ngừng 24--25/12/2026
\item Lộ trình đã \al{lùi hai lần}; UPCoM không thuộc đợt này
\end{itemize}
\end{frame}

\begin{frame}
\frametitle{Hà Nội cần biết mất gì, được gì khi không còn sàn cổ phiếu niêm yết}
\begin{columns}[T]
\column{0.50\textwidth}
\textsb{Bài toán số 19 của Hà Nội}\\[1mm]
{\small ``Phương thức nào để xây dựng Hà Nội thành trung tâm tài chính và dịch vụ chất lượng cao quốc tế?''}\\[1mm]
{\footnotesize QĐ 1193/QĐ-UBND, 2026}
\vspace{4mm}
\begin{itemize}
\item NQ 222/2025/QH15: trung tâm tài chính quốc tế tại TP.HCM và Đà Nẵng
\item Chưa có số đo về chất lượng thị trường của Hà Nội
\end{itemize}
\column{0.45\textwidth}
\textsb{GFCI 40 (9/2026)}\\[2mm]
\begin{tabular}{@{}lr@{}}
\toprule
Trung tâm & Hạng / 117\\
\midrule
TP. Hồ Chí Minh & 67\\
Đà Nẵng & 71\\
Hà Nội & \al{không có}\\
\bottomrule
\end{tabular}\\[1mm]
{\footnotesize Wardle \& Mainelli (2026)}
\end{columns}
\end{frame}

\begin{frame}
\frametitle{Ba câu hỏi, bốn giả thuyết; trọng tâm là tác động của việc chuyển sàn}
\begin{center}
{\small
\begin{tabularx}{\textwidth}{@{}>{\raggedright\arraybackslash}p{0.33\textwidth} L l@{}}
\toprule
Câu hỏi & Giả thuyết & Vai trò\\
\midrule
1. Thị trường cổ phiếu phát triển có đi cùng điểm GFCI cao hơn? & \HH{1}: liên hệ dương, sau khi kiểm soát thể chế, thu nhập & bối cảnh\\
2. Trước hợp nhất, HNX kém HOSE bao nhiêu? & \HH{2}: chỉ số chất lượng HNX thấp hơn HOSE (6/2009--6/2025) & thước đo nền\\
3. Chuyển sàn thay đổi thanh khoản và khả năng đầu tư bao nhiêu? & \HH{3}: độ phi thanh khoản giảm so với đối chứng & \al{trọng tâm}\\
 & \HH{4}: khoảng cách khả năng đầu tư thu hẹp & ước lượng\\
\bottomrule
\end{tabularx}}
\end{center}
\vspace{1mm}
{\small \HH{1}--\HH{3}: kiểm định xác nhận, hiệu chỉnh Holm 5\%. \HH{4}: chỉ ước lượng (ít cửa sổ).}
\end{frame}

\begin{frame}
\frametitle{Đóng góp: đo vai trò của sàn ở một thị trường mới nổi, bằng một sự kiện ngoại sinh}
\begin{center}
{\small
\begin{tabularx}{\textwidth}{@{}>{\raggedright\arraybackslash}p{0.27\textwidth} L L@{}}
\toprule
Nhóm tài liệu & Đã có & Đề tài bổ sung\\
\midrule
Trung tâm tài chính & Naeem (2023): 17 yếu tố, không có thị trường cổ phiếu & Đưa riêng thị trường cổ phiếu vào mô hình GFCI\\
Hợp nhất sàn & Euronext, Tokyo--Osaka: thị trường phát triển, chỉ đo thanh khoản & Thị trường mới nổi; suy luận cho \al{một sự kiện}; thêm khả năng đầu tư\\
Việt Nam & Tập trung vào HOSE & So sánh HNX với HOSE; đo đợt hợp nhất 2026\\
\bottomrule
\end{tabularx}}\\[1mm]
{\footnotesize Trong phạm vi tìm kiếm đến 30/9/2026.}
\end{center}
\end{frame}

\begin{frame}
\frametitle{Khung phân tích: từ chất lượng trước hợp nhất đến hàm ý cho Hà Nội}
\begin{center}
\begin{tikzpicture}[node distance=5mm and 10mm, every node/.append style={inner sep=4pt}]
\node[box, text width=3.3cm] (a) {\textsb{Trước hợp nhất}\\HNX so với HOSE (\HH{2})\\6/2009--6/2025};
\node[box, text width=3.3cm, right=of a, draw=RubineRed, very thick] (b) {\textsb{Sự kiện 28/12/2026}\\Thanh khoản (\HH{3})\\Khả năng đầu tư (\HH{4})};
\node[box, text width=3.3cm, right=of b] (c) {\textsb{Hàm ý cho Hà Nội}\\Chuyên môn hóa thị trường\\Hạ tầng dữ liệu};
\node[box, text width=3.3cm, below=of b] (d) {\textsb{Nhóm đối chứng}\\Cổ phiếu HOSE đối sánh\\Cổ phiếu UPCoM};
\node[box, text width=3.3cm, below=of c] (e) {\textsb{Bằng chứng quốc tế}\\Điểm GFCI (\HH{1})\\Thể chế, thu nhập};
\draw[arr] (a) -- (b); \draw[darr] (b) -- (c);
\draw[arr] (d) -- node[right, font=\footnotesize]{khác biệt kép} (b);
\draw[darr] (e) -- node[right, font=\footnotesize]{bối cảnh} (c);
\end{tikzpicture}
\end{center}
{\footnotesize Nét liền: kiểm định hoặc ước lượng. Nét đứt: tổng hợp định tính.}
\end{frame}

%=================================================================== DỮ LIỆU, ĐO LƯỜNG
\begin{frame}
\heading{Dữ liệu và đo lường}
\end{frame}

\begin{frame}
\frametitle{Ba giai đoạn: tham chiếu, chuyển tiếp bị loại, 24 tháng sau sự kiện}
\begin{center}
\begin{tikzpicture}[x=0.62cm, y=0.9cm]
\draw[arr] (0,0) -- (21,0) node[right, font=\small]{năm};
\foreach \y/\l in {0/2009,4/2013,8/2017,12/2021,16/2025,20/2029}{\draw (\y,0.1)--(\y,-0.1) node[below, font=\footnotesize]{\l};}
\draw[thick, fill=LightGray] (0.5,0.45) rectangle (16.5,1.25); \node[font=\small] at (8.5,0.85){Tham chiếu 6/2009--6/2025};
\draw[thick, dashed] (16.5,0.45) rectangle (18,1.25);
\draw[thick, fill=RubineRed!15] (18,0.45) rectangle (20,1.25);
\draw[very thick] (18,-0.15) -- (18,1.6);
\node[anchor=south east, font=\footnotesize, align=right] at (17.2,1.7){chuyển tiếp\\7/2025--12/2026};
\node[anchor=south west, font=\footnotesize, align=left] at (18.8,1.7){sau sự kiện\\1/2027--12/2028};
\end{tikzpicture}
\end{center}
\begin{itemize}
\item Dữ liệu ngày của sàn niêm yết HNX, UPCoM, HOSE, \al{kể cả mã đã hủy niêm yết} (tránh thiên lệch sống sót)
\item Chuyển tiếp bị loại khỏi kiểm định chính: thị trường đã biết trước sự kiện từ 1/7/2025
\item Quốc tế: điểm GFCI kỳ 1--40, chỉ báo thị trường cổ phiếu, quản trị nhà nước
\end{itemize}
\end{frame}

\begin{frame}
\frametitle{Chỉ số chất lượng thị trường vốn: ba chiều, sáu chỉ báo, trọng số bằng nhau}
\[
\MQI_{jt} = \frac{\vn{TK}_{jt} + \vn{HQ}_{jt} + \vn{MB}_{jt}}{3}
\]
\begin{center}
{\small
\begin{tabularx}{0.92\textwidth}{@{}l L@{}}
\toprule
Chiều & Chỉ báo\\
\midrule
Thanh khoản (TK) & độ phi thanh khoản Amihud; tỷ lệ ngày lợi suất 0; chênh lệch Corwin--Schultz\\
Hiệu quả thông tin (HQ) & độ lệch tỷ số phương sai khỏi 1; độ trễ giá\\
Minh bạch (MB) & số ngày trễ hạn công bố báo cáo quý\\
\bottomrule
\end{tabularx}}
\end{center}
\begin{itemize}
\item Mỗi chỉ báo đổi thành điểm z theo mẫu gộp các sàn trong giai đoạn tham chiếu
\item \HH{2}: chênh lệch HOSE -- HNX theo tháng; giá trị p bằng bootstrap khối dừng
\end{itemize}
\end{frame}

\begin{frame}
\frametitle{Thách thức: biên độ và bước giá đổi cùng lúc với sàn}
\begin{columns}[T]
\column{0.48\textwidth}
\textsb{Quy tắc thay đổi khi chuyển sàn}\\[2mm]
{\small
\begin{tabular}{@{}lcc@{}}
\toprule
 & HNX & HOSE\\
\midrule
Biên độ & ±10\% & ±7\%\\
Bước giá & 100 đ & 10 / 50 / 100 đ\\
\bottomrule
\end{tabular}}\\[3mm]
{\small Biên độ hẹp $\Rightarrow$ lợi suất bị cắt $\Rightarrow$ Amihud \al{giảm cơ học}}
\column{0.48\textwidth}
\textsb{Cách xử lý}
\begin{itemize}
\item Tính Amihud trên ngày có $|r|$ \al{dưới một ngưỡng chung}, bỏ ngày liền sau
\item Ngưỡng dưới mức tăng tối đa ở ±7\%: giá 1.140 đ chỉ tăng được \al{6,14\%}
\item Loại mã giá quá thấp; báo cáo theo nhóm giá
\end{itemize}
\end{columns}
\end{frame}

%=================================================================== NHẬN DẠNG
\begin{frame}
\heading{Thiết kế nhận dạng}
\end{frame}

\begin{frame}
\frametitle{Khác biệt kép: so cổ phiếu chuyển sàn với cổ phiếu HOSE đối sánh}
\[
A_{it} = \alpha_i + \tau_t + \delta D_i\Post_t + \varepsilon_{it}
\]
\begin{columns}[T]
\column{0.50\textwidth}
\begin{itemize}
\item $A_{it}$: log Amihud (ngày dưới ngưỡng chung)
\item $D_i = 1$: mã chuyển sàn; $\Post_t = 1$ từ 1/2027
\item $\alpha_i,\ \tau_t$: hiệu ứng cố định mã, tháng
\item \HH{3}: $\delta < 0$
\end{itemize}
\column{0.46\textwidth}
\textsb{Ví dụ minh họa}\\[1mm]
{\small
\begin{tabular}{@{}lccc@{}}
\toprule
 & Trước & Sau & Đổi\\
\midrule
Chuyển sàn & 1,80 & 1,40 & $-0{,}40$\\
Đối chứng & 1,20 & 1,05 & $-0{,}15$\\
\midrule
$\delta$ & & & \al{$-0{,}25$}\\
\bottomrule
\end{tabular}}\\[1mm]
{\small $e^{-0{,}25} - 1 \approx -22\%$}
\end{columns}
\end{frame}

\begin{frame}
\frametitle{Nhóm đối chứng và các cú sốc trùng thời gian được xử lý riêng}
\begin{columns}[T]
\column{0.42\textwidth}
\textsb{Đối chứng}
\begin{itemize}
\item Chính: mã HOSE ghép theo vốn hóa, ngành, Amihud tại 30/6/2025
\item Phụ: mã UPCoM
\item Kiểm tra xu hướng song song bằng nghiên cứu sự kiện; độ nhạy theo Rambachan \& Roth (2023)
\end{itemize}
\column{0.54\textwidth}
\textsb{Cú sốc trùng thời gian}\\[1mm]
{\small
\begin{tabularx}{\linewidth}{@{}L L@{}}
\toprule
Cú sốc & Xử lý\\
\midrule
FTSE Russell (9/2026--9/2027) & Loại mã đủ điều kiện khỏi đối chứng\\
Phân bảng HOSE & Thêm biến giả nhóm bảng\\
VN-Index nhận mã mới & Đo mức chịu cú sốc chỉ số\\
Nới biên độ, xuyên trưa & Cắt cửa sổ sau (phụ)\\
\bottomrule
\end{tabularx}}
\end{columns}
\end{frame}

\begin{frame}
\frametitle{Chỉ có một sự kiện, nên suy luận bằng bootstrap tự hồi quy}
\begin{columns}[T]
\column{0.48\textwidth}
\textsb{Vấn đề}
\begin{itemize}
\item 299 mã nhưng chỉ \al{một lần} chuyển sàn
\item Coi mã là độc lập $\Rightarrow$ sai số chuẩn quá nhỏ\\{\footnotesize (Conley \& Taber, 2011)}
\end{itemize}
\textsb{Cách làm}
\begin{itemize}
\item Chuỗi $d_t$: chênh lệch log Amihud, chuyển sàn trừ đối chứng
\item Thống kê: trung bình $d_t$ sau trừ trước
\end{itemize}
\column{0.48\textwidth}
\begin{center}
\begin{tikzpicture}[node distance=3mm, every node/.append style={font=\footnotesize, inner sep=3pt}]
\node[box, text width=5cm] (s1) {Ước lượng AR trên giai đoạn tham chiếu};
\node[box, text width=5cm, below=of s1] (s2) {Tạo chuỗi giả theo mô hình AR, rút phần dư có hoàn lại};
\node[box, text width=5cm, below=of s2] (s3) {Tính thống kê trên mỗi chuỗi giả};
\node[box, text width=5cm, below=of s3, draw=RubineRed, very thick] (s4) {Giá trị p = tỷ lệ chuỗi giả cực đoan hơn};
\draw[arr] (s1)--(s2); \draw[arr] (s2)--(s3); \draw[arr] (s3)--(s4);
\end{tikzpicture}
\end{center}
\end{columns}
\end{frame}

\begin{frame}
\frametitle{Lực kiểm định phụ thuộc vào số tháng, không phụ thuộc vào số mã}
\begin{columns}[T]
\column{0.50\textwidth}
\vspace{2mm}
\textsb{Hiệu ứng nhỏ nhất phát hiện được}\\[1mm]
{\small (lực 80\%, mức một phía 0,05/3; $\sigma = 0{,}2$ minh họa)}\\[2mm]
{\small
\begin{tabular}{@{}cccc@{}}
\toprule
$\phi$ & Theo $\sigma$ & Log Amihud & Amihud giảm\\
\midrule
0 & 0,64 & 0,13 & 12\%\\
0,6 & 1,24 & 0,25 & \al{22\%}\\
0,8 & 1,76 & 0,35 & 30\%\\
\bottomrule
\end{tabular}}
\column{0.46\textwidth}
\begin{itemize}
\item $\phi$: tự tương quan của $d_t$; chuỗi càng ``dính'', cần tác động càng lớn
\item Cửa sổ trước rút còn 24 tháng $\Rightarrow$ ngưỡng tăng khoảng 1,3 lần
\item Không bác bỏ $\neq$ không có tác động: báo cáo khoảng tin cậy
\item $\sigma$, $\phi$ thật ước lượng lại, ghi vào tệp đăng ký
\end{itemize}
\end{columns}
\end{frame}

\begin{frame}
\frametitle{Tác động được tách theo ba kênh và theo trụ sở doanh nghiệp}
\vspace{2mm}
\begin{center}
{\small
\begin{tabularx}{\textwidth}{@{}>{\raggedright\arraybackslash}p{0.22\textwidth} L L@{}}
\toprule
Thay đổi khi chuyển sàn & Mã chịu mạnh nhất & Thước đo (dữ liệu đến 30/6/2025)\\
\midrule
Biên độ ±10\% $\rightarrow$ ±7\% & Mã hay tăng giảm 7--10\% & Tỷ lệ ngày $|r|$ từ 7\% đến dưới 10\%\\
Bước giá nhỏ đi & Mã giá thấp & Nghịch đảo giá đóng cửa bình quân\\
Vào chỉ số HOSE & Mã vốn hóa lớn so với giao dịch & Vốn hóa / giá trị giao dịch bình quân ngày\\
\bottomrule
\end{tabularx}}
\end{center}
\begin{itemize}
\item Tương tác $D_i\Post_t E_{si}$; suy luận cùng cách với $\delta$; là phân tích phụ, ngoài thủ tục Holm
\item Trụ sở tại Hà Nội: kiểm tra kênh chú ý và thông tin gắn với nơi đặt sàn
\end{itemize}
\end{frame}

%=================================================================== KHẢ NĂNG ĐẦU TƯ, H1
\begin{frame}
\heading{Khả năng đầu tư và bằng chứng quốc tế}
\end{frame}

\begin{frame}
\frametitle{Khả năng đầu tư: nhà đầu tư thiệt bao nhiêu với rổ cổ phiếu chuyển sàn}
\begin{columns}[T]
\column{0.50\textwidth}
\vspace{-2mm}
\begin{align*}
U &= \mu - \frac{\rho}{2}\sigma^2\\
G_m &= U_m(\vn{H}) - U_m(\vn{CS})\\
\bar{G} &= \frac{G_1 + G_2 + \dots + G_7}{7}
\end{align*}
{\small $U$: lợi suất tương đương chắc chắn sau chi phí; rổ CS: tối đa 30 mã chuyển sàn; rổ H: mã HOSE ghép cặp}
\column{0.46\textwidth}
\begin{itemize}
\item \al{7 mô hình phân bổ} \MM{1}--\MM{7}: từ 1/N đến Black-Litterman chuyển chế độ
\item Bình quân để không phụ thuộc một mô hình
\item 20 cửa sổ trước, 8 cửa sổ sau
\item \HH{4}: $\bar{G}$ thu hẹp sau sự kiện (ước lượng)
\end{itemize}
\end{columns}
\end{frame}

\begin{frame}
\frametitle{Bằng chứng quốc tế: thị trường cổ phiếu và điểm GFCI là liên hệ có điều kiện}
\[
Z_{it} = \alpha_i + \tau_t + \beta \SMD_{it} + \lambda W_{it} + \varepsilon_{it}
\]
\begin{columns}[T]
\column{0.50\textwidth}
\begin{itemize}
\item $Z_{it}$: điểm GFCI chuẩn hóa theo kỳ
\item $\SMD_{it}$: vòng quay, vốn hóa / GDP
\item $W_{it}$: thu nhập, quản trị nhà nước
\item \HH{1}: $\beta > 0$; sai số phân cụm theo quốc gia
\end{itemize}
\column{0.46\textwidth}
\textsb{Giới hạn nêu thẳng}
\begin{itemize}
\item GFCI đã chứa dữ liệu giao dịch cổ phiếu
\item $\SMD$ ít đổi theo thời gian
\item $\Rightarrow$ không gọi là tác động nhân quả
\end{itemize}
\end{columns}
\end{frame}

%=================================================================== KẾ HOẠCH, SẢN PHẨM
\begin{frame}
\heading{Kế hoạch và sản phẩm}
\end{frame}

\begin{frame}
\frametitle{Kế hoạch phân tích được khóa trước ngày 28/12/2026}
\begin{columns}[T]
\column{0.46\textwidth}
\textsb{Đăng ký trước trên OSF Registries}
\begin{itemize}
\item Chỉ báo, ngưỡng, quy tắc ghép
\item Đặc tả phương trình, suy luận
\item Mức ý nghĩa, hiệu chỉnh Holm
\item Mọi thay đổi sau đó ghi là lệch kế hoạch
\end{itemize}
\column{0.50\textwidth}
\textsb{Phương án dự phòng}\\[1mm]
{\small
\begin{tabularx}{\linewidth}{@{}L L@{}}
\toprule
Rủi ro & Xử lý\\
\midrule
Chuyển sàn lùi, nhiều đợt & Callaway \& Sant'Anna (2021)\\
UPCoM chuyển sớm & Dựa vào đối chứng HOSE\\
Không song song & Đối chứng tổng hợp\\
Không dùng được GFCI & \HH{1} ở dạng mô tả\\
\bottomrule
\end{tabularx}}
\end{columns}
\end{frame}

\begin{frame}
\frametitle{Bốn năm, sáu chỉ tiêu kiểm chứng được}
\begin{columns}[T]
\column{0.52\textwidth}
{\small
\begin{tabularx}{\linewidth}{@{}l L@{}}
\toprule
Năm & Việc chính\\
\midrule
1 & Đăng ký trước; dữ liệu; chỉ số chất lượng; \HH{2}\\
2 & Bài báo 1; \HH{1}; kết quả giữa kỳ \HH{3}\\
3 & 24 tháng sau sự kiện; \HH{3}, \HH{4}; bài báo 2, 3\\
4 & Báo cáo định hướng cho Hà Nội; luận án\\
\bottomrule
\end{tabularx}}
\column{0.44\textwidth}
\textsb{Sản phẩm}
\begin{itemize}
\item Tệp đăng ký trước có dấu thời gian
\item Chỉ số chất lượng thị trường và chỉ số khả năng đầu tư, mã nguồn mở
\item Ba bài báo, ít nhất một bài WoS Q1/Q2
\item Báo cáo cho Sở KH\&CN Hà Nội, ưu tiên HNX dùng
\end{itemize}
\end{columns}
\end{frame}

%=================================================================== KẾT
\begin{frame}
\frametitle{Ba điều cần nhớ}
\vspace{4mm}
\begin{enumerate}
\item \textsb{Cơ hội đo hiếm có:} 299 mã đổi sàn cùng ngày, không do doanh nghiệp chọn
\item \textsb{Thiết kế chặt:} khác biệt kép, ngưỡng lợi suất chung, bootstrap cho một sự kiện, đăng ký trước
\item \textsb{Sản phẩm dùng được dù kết quả ra sao:} hai chỉ số và báo cáo định hướng cho thị trường vốn Hà Nội
\end{enumerate}
\vspace{6mm}
{\small Câu hỏi còn mở: Hà Nội nên chuyên môn hóa vào trái phiếu, phái sinh hay UPCoM?}
\end{frame}

\begin{frame}
\frametitle{Tài liệu chính}
{\small
\begin{itemize}
\item Amihud, Y. (2002). Illiquidity and stock returns. \emph{Journal of Financial Markets, 5}(1), 31--56.
\item Conley, T. G., \& Taber, C. R. (2011). Inference with ``difference in differences'' with a small number of policy changes. \emph{Review of Economics and Statistics, 93}(1), 113--125.
\item Hisada, T. (2018). The effect of investor sentiment toward an exchange merger on liquidity. \emph{International Journal of Economics and Financial Issues, 8}(2), 315--318.
\item Naeem, M. M. F. (2023). What makes an International Financial Centre (IFC) competitive? [Luận án tiến sĩ, Durham University].
\item Nielsson, U. (2009). Stock exchange merger and liquidity: The case of Euronext. \emph{Journal of Financial Markets, 12}(2), 229--267.
\item Rambachan, A., \& Roth, J. (2023). A more credible approach to parallel trends. \emph{Review of Economic Studies, 90}(5), 2555--2591.
\end{itemize}}
\end{frame}

\begin{frame}
\begin{center}
{\LARGE\sbseries Xin cảm ơn}\\[4mm]
Nguyễn Văn Trung
\end{center}
\end{frame}

\lastslide

%=================================================================== PHỤ LỤC
\appendix
\begin{frame}
\begin{center}{\Large\sbseries Slide phụ lục}\end{center}
\end{frame}

\begin{frame}
\frametitle{Phụ lục: bảy mô hình phân bổ}
{\small
\begin{tabularx}{\textwidth}{@{}l L l@{}}
\toprule
Mã & Mô hình & Nguồn\\
\midrule
\MM{1} & Phân bổ đều 1/N & DeMiguel và cộng sự (2009)\\
\MM{2} & Phương sai nhỏ nhất, hiệp phương sai co rút & Ledoit \& Wolf (2004)\\
\MM{3} & Trung bình-phương sai, ước lượng mẫu & Markowitz (1952)\\
\MM{4} & Black-Litterman & Black \& Litterman (1992)\\
\MM{5} & Black-Litterman có chuyển chế độ & Nguyễn \& Nguyễn (2026a)\\
\MM{6} & Danh mục vững chắc, nhiều niềm tin tiên nghiệm & Garlappi và cộng sự (2007)\\
\MM{7} & Giảm rủi ro theo xác suất trạng thái căng thẳng & Kritzman và cộng sự (2012)\\
\bottomrule
\end{tabularx}}\\[2mm]
{\small Danh mục chỉ mua; ước lượng lại hằng tháng trên cửa sổ 104 tuần; báo cáo cả khi bỏ \MM{5}.}
\end{frame}

\begin{frame}
\frametitle{Phụ lục: nghiên cứu sự kiện kiểm tra xu hướng song song}
\[
A_{it} = \alpha_i + \tau_t + \sum_k \delta_k D_i T_{kt} + \varepsilon_{it}
\]
\begin{itemize}
\item $k$ từ 1/2024 đến 12/2028; tháng gốc 6/2025; trước 1/2024 gộp một nhóm
\item Kiểm định đồng thời $\delta_k = 0$ từ 1/2024 đến 5/2025
\item Độ nhạy trước vi phạm xu hướng song song: Rambachan \& Roth (2023)
\item 18 tháng chuyển tiếp giữa tháng gốc và giai đoạn sau $\Rightarrow$ khoảng tin cậy có thể rộng
\end{itemize}
\end{frame}

\begin{frame}
\frametitle{Phụ lục: vì sao không dùng sai số chuẩn phân cụm theo mã}
\begin{itemize}
\item Mọi mã cùng chịu một sự kiện, cùng một ngày
\item Cú sốc tác động như nhau lên mọi mã: đã nằm trong $\tau_t$
\item Cú sốc tác động khác nhau lên hai nhóm: làm sai số của các mã chuyển sàn tương quan trong cùng tháng
\item Thêm mã không làm giảm loại cú sốc này $\Rightarrow$ thông tin nằm ở số tháng
\item Kiểm tra bổ sung: ngày sự kiện giả trong giai đoạn tham chiếu (Conley \& Taber, 2011); chỉ khoảng sáu cửa sổ độc lập nên không dùng làm suy luận chính
\end{itemize}
\end{frame}

\begin{frame}
\frametitle{Phụ lục: hiệu chỉnh Holm cho \HH{1}--\HH{3}}
\begin{columns}[T]
\column{0.48\textwidth}
\begin{itemize}
\item Sắp giá trị p từ nhỏ đến lớn
\item So lần lượt với 0,05/3; 0,05/2; 0,05
\item Dừng ở giá trị đầu tiên không đạt
\item \HH{4} ngoài họ: chỉ 28 cửa sổ
\end{itemize}
\column{0.48\textwidth}
\textsb{Ví dụ minh họa}\\[1mm]
{\small
\begin{tabular}{@{}lccl@{}}
\toprule
 & p & Ngưỡng & Kết quả\\
\midrule
\HH{2} & 0,004 & 0,0167 & đạt\\
\HH{3} & 0,020 & 0,025 & đạt\\
\HH{1} & 0,030 & 0,050 & đạt\\
\bottomrule
\end{tabular}}
\end{columns}
\end{frame}

\begin{frame}
\frametitle{Phụ lục: mốc pháp lý}
{\small
\begin{tabularx}{\textwidth}{@{}l L@{}}
\toprule
Thời điểm & Sự kiện\\
\midrule
2016 & Quy chế giao dịch HOSE (QĐ 341) và HNX (QĐ 529): bước giá\\
2021--2025 & TT 57/2021, 69/2023, 139/2025: hạn chuyển sàn lùi đến 31/12/2026\\
1/7/2025 & HNX ngừng nhận hồ sơ niêm yết mới\\
9/2026--9/2027 & FTSE Russell đưa cổ phiếu Việt Nam vào chỉ số theo bốn đợt\\
30/9/2026 & VNX công bố lịch chuyển sàn\\
28/12/2026 & Phiên đầu tại HOSE\\
1/2027--12/2028 & 24 tháng sau sự kiện\\
\bottomrule
\end{tabularx}}
\end{frame}

\end{document}
